\documentclass[11pt]{article}
\usepackage[a4paper,margin=22mm]{geometry}
\usepackage[no-math]{fontspec}
\setmainfont{Latin Modern Roman}
\usepackage{amsmath,amssymb,amsthm,xfrac,xspace,hyperref,graphicx,comment}
\excludecomment{DefTralics}
\providecommand{\C}{}
\providecommand{\bibliofont}{\small}
\newcommand{\ie}{i.e.,\xspace}
\newcommand{\eg}{e.g.,\xspace}
\newcommand{\etc}{etc.\xspace}
\newcommand{\cf}{cf.\xspace}
\newcommand{\resp}{resp.\xspace}
\newcommand{\smaller}{\small}
\newcommand{\Subsection}{\subsection}
\newcommand{\Subsubsection}{\subsubsection}
\newcommand{\skpt}{\smallskip}
\emergencystretch=3em

% Begin literal retained input author-preamble.tex
\newenvironment{enumeratei}
{\bgroup\def\theenumi{\roman{enumi}}\def\theenumii{\arabic{enumii}}\begin{enumerate}}
{\end{enumerate}\egroup}
\newenvironment{enumerateA}
{\bgroup\def\theenumi{\Alph{enumi}}\begin{enumerate}}
{\end{enumerate}\egroup}
\newcommand{\sbullet}{{\scriptscriptstyle\bullet}}
\newcommand\mto{\mathchoice{\longmapsto}{\mapsto}{\mapsto}{\mapsto}}
\newcommand\implique{\mathchoice{\implies}{\Rightarrow}{\Rightarrow}{\Rightarrow}}
\newcommand\ssi{\mathchoice{\iff}{\Leftrightarrow}{\Leftrightarrow}{\Leftrightarrow}}
\newcommand{\rb}{\mathrm{b}}
\newcommand{\ru}{\mathrm{ru}}
\newcommand{\Changel}{\mathcode`l="7160}
\newcommand{\Changelback}{\mathcode`l="716C}
\newcommand{\NoChangel}[1]{%
\expandafter\let\csname old\string#1\endcsname=#1
\let#1=\relax
\newcommand{#1}{\mathcode`l="716C\csname old\string#1\endcsname\mathcode`l="7160 }%
}

\NoChangel{\log}
\NoChangel{\ln}
\NoChangel{\lim}
\NoChangel{\limsup}
\NoChangel{\liminf}
\NoChangel{\varlimsup}
\NoChangel{\varliminf}
\NoChangel{\varinjlim}
\NoChangel{\varprojlim}

\usepackage{mathrsfs}
\let\mathcal\mathscr

\usepackage{pgf,tikz}
\usetikzlibrary{arrows}

\theoremstyle{plain}
\newtheorem{thm}{Theorem}[section]
\newtheorem{lem}[thm]{Lemma}
\newtheorem{prop}[thm]{Proposition}
\newtheorem{cor}[thm]{Corollary}
\newtheorem*{claim*}{Claim}

\theoremstyle{definition}
\newtheorem{defn}[thm]{Definition}
\newtheorem{example}[thm]{Example}
\newtheorem{rem}[thm]{Remark}

\newcommand*{\mydprime}{^{\prime\prime}\mkern-1.2mu}

\newcommand{\bA}{{\boldsymbol{A}}}
\newcommand{\bB}{{\boldsymbol{B}}}
\newcommand{\bC}{{\boldsymbol{C}}}
\newcommand{\bD}{{\boldsymbol{D}}}
\newcommand{\bE}{{\boldsymbol{E}}}
\newcommand{\bF}{{\boldsymbol{F}}}
\newcommand{\bK}{{\boldsymbol{K}}}
\newcommand{\bS}{{\boldsymbol{S}}}
\newcommand{\bCT}{{\boldsymbol{CT}}}

\newcommand{\CC}{\mathbf{C}}
\newcommand{\NN}{\mathbf{N}}
\newcommand{\ZZ}{\mathbf{Z}}
\newcommand{\RR}{\mathbf{R}}
\newcommand{\QQ}{\mathbf{Q}}

\newcommand{\sU}{\mathscr{U}}

\newcommand{\cC}{\mathcal{C}}
\newcommand{\cCT}{\mathcal{CT}}
\newcommand{\cE}{\mathcal{E}}
\newcommand{\cK}{\mathcal{K}}
\newcommand{\cL}{\mathcal{L}}
\newcommand{\cS}{\mathcal{S}}
\newcommand{\cT}{\mathcal{T}}
\newcommand{\cU}{\mathcal{U}}
\newcommand{\cV}{\mathcal{V}}

\newcommand{\sinf}{\cS_\infty}

\newcommand{\sfc}{\mathsf{c}}

\let\se\subseteq
\let\action\curvearrowright

\newcommand{\da}[1]{{#1\!\downarrow}}

\DeclareMathOperator{\supp}{Supp}
\DeclareMathOperator{\Aut}{Aut}
\DeclareMathOperator{\Int}{Int}
\DeclareMathOperator{\Homeo}{Homeo}
\DeclareMathOperator{\Ends}{Ends}
\DeclareMathOperator{\Sym}{Sym}
\DeclareMathOperator{\Br}{Br}
\DeclareMathOperator{\Fix}{Fix}
\DeclareMathOperator{\Stab}{Stab}
\DeclareMathOperator{\CLO}{CCLO}
\DeclareMathOperator{\LO}{LO}


% End literal retained input author-preamble.tex


\begin{document}
\begin{center}{\Large Stable item 1927}\end{center}
\paragraph{Source and attribution.}
Bruno Duchesne, \emph{Topological properties of Ważewski dendrite groups}, Journal de l’École polytechnique — Mathématiques 7 (2020), publisher version, DOI 10.5802/jep.121. \href{https://jep.centre-mersenne.org/articles/10.5802/jep.121/}{Publisher article}. Authors' text: \href{https://creativecommons.org/licenses/by/4.0/}{CC BY 4.0}.
\paragraph{Editorial problem boundary.}
Prove only the original Theorem 6.2: if a closed subgroup \(G\) of the full symmetric group on a countable set has a comeager conjugacy class and the small index property, then its unrestricted permutational wreath product \(G^\Omega\rtimes\mathrm{Sym}(\Omega)\), with the printed imprimitive Polish topology, has the small index property. Exact definitions, full Theorem 6.2, Lemma 6.4 with its proof, and the entire wreath-product proof are retained. The dendrite application is not selected.
\paragraph{Difficulty (editorial): H2.}
A small-index subgroup is forced to contain a full supported wreath factor by an infinite disjoint-moiety projection and comeager conjugacy argument. Almost-disjoint continuum families expand it to a cofinite supported factor, leaving finitely many base-group subgroups whose openness closes the proof.
\paragraph{Editorial transcription note.}
Original author language, mathematical content and ordering are unchanged. Publisher front matter is replaced by this independent layout wrapper. Original native statement, proof and macros remain editable; bibliography is recovered from the matching cached publisher HTML. No assistant-generated proof or mathematical repair is supplied.
\clearpage
\setcounter{section}{5}

% Begin literal retained input author-section.tex
\section{Small index subgroups}
Let us recall that a Polish group has the \emph{small index property} if any subgroup of small index, \ie of index less than $2^{\aleph_0}$, is open. For example $\sinf$ and $\Aut(\QQ,<)$ have this property.

% End literal retained input author-section.tex


% Begin literal retained input author-definitions.tex
Let $\Omega$ be a countable infinite set with full permutation group $\sinf$. For a group $G$, we denote by $G\wr \sinf$ the (unrestricted permutational) wreath product $G^\Omega\rtimes\sinf$. The action of $\sinf$ on $G^\Omega$ is by permutation of the coordinates. If $\sigma\in\sinf$ and $(g_\omega)_\omega\in G^\Omega$ then
\[
\sigma\cdot(g_\omega)=(g_{\sigma^{-1}\omega})_\omega.
\]
If the role of $\Omega$ has to be emphasized, we denote the above wreath product $G\wr_\Omega \sinf$.

In the particular case where $G$ is a closed subgroup of $\sinf$ acting on a countable set $\Lambda$ and another copy of $\sinf$ acts as above on the countable set $\Omega$, the wreath product $G\wr\sinf$ acts on $\Lambda\times \Omega$ with the \emph{imprimitive action}. This action is given by the following formula:
\[
\left((g_\omega),\sigma\right)\cdot(\lambda,\omega')=\bigl(g_{\sigma(\omega')}\lambda,\sigma(\omega')\bigr).
\]
This action embeds $G\wr\sinf$ as a closed subgroup of the symmetric group of $\Lambda\times \Omega$ and thus it has a natural Polish topology and we will consider this group with this topology.

% End literal retained input author-definitions.tex

\setcounter{thm}{1}
% Begin literal retained input author-statement.tex
\begin{thm}\label{wr} Let $G$ be a closed subgroup of $\sinf$ with a comeager conjugacy class and the small index property. The wreath product $W=G\wr\sinf$ is a Polish group with the small index property.\end{thm}

% End literal retained input author-statement.tex

\setcounter{thm}{3}
% Begin literal retained input author-proof.tex
\begin{lem}\label{nsi} Let $G$ be a Polish group with a comeager conjugacy class. If $N$ is a normal subgroup of small index then $N=G$.
\end{lem}

\begin{proof}Let $\cC$ be the comeager conjugacy class. It suffices to show that $\cC\cap N\neq\emptyset$. Otherwise, by normality, $\cC\subset N$ and thus $G=\cC\cdot\cC\subset N$.

Since $N$ has small index then $N$ is not meager \cite[Th.\,4.1]{MR1231710} (see also \cite[Lem.\,6.8]{Kechris-Rosendal} for a very short proof), so $N\cap\cC\neq\emptyset$ and we are done.\end{proof}

Our proof that $G\wr\sinf$ has the small index subgroup property borrows the original ideas that lead to prove the property for $\sinf$ \cite[Th.\,1]{MR859950}. We not only use the result but also the proof itself and thus reproduce some of the arguments there. Let us recall that a moiety of $\Omega$ is a subset $\Sigma$ that is infinite and co-infinite.

\begin{proof}[Proof of Theorem~\ref{wr}] Let $H$ be a subgroup of $W$ of small index. Let $(\Sigma_i)$ be an infinite collection of disjoint moieties of $\Omega$. Let $W_i=G\wr_{\Sigma_i} \Sym(\Sigma_i)\leq\Sym(\Lambda\times \Omega)$ be the subgroup of permutations supported on $\Lambda\times\Sigma_i$. More precisely, $\left((g_\omega),\sigma\right)\in W_i$ if $\supp(\sigma)\subset \Sigma_i$ and $g_\omega=e$ for $\omega\notin \Sigma_i$. By disjointness of the supports, for $i\neq j$, $W_i\cap W_j$ is trivial and these two subgroups commute. Let $P$ be the subgroup $\Sym(\Lambda\times \Omega)$ isomorphic to $\prod_i W_i$. Let $H_i$ be the projection of $H\cap P$ on $W_i$. We have
\[
\prod_{i}|W_i\colon H_i|=\Bigl|P\colon \prod_{i}H_i\Bigr|\leq|P\colon H\cap P|\leq|W\colon H|<2^{\aleph_0}.
\]

This implies that for all but finitely many, $W_i=H_i$. We fix such an $i$ and simply note $\Sigma=\Sigma_i$ and $W'=W_i$. We denote by $G^\Sigma$ the subgroup of $G\wr\sinf$ with elements of the form $((g_\omega),e)$ and $g_\omega=e$ for $\omega\notin \Sigma$. Let $g\in H\cap G^\Sigma$ and $g'\in G^\Sigma$. There exists $h\in H\cap P$ such that $\pi_i(h)=g'$, where $\pi_i\colon P\to W_i$ is the projection. One has $g'g(g')^{-1}= hgh^{-1}\in H\cap G^\Sigma$. So, the subgroup $H\cap G^\Sigma$ is a normal subgroup of $G^\Sigma$ of small index. The group $G^\Sigma$ has a comeager conjugacy class because of \cite[Lem.\,11]{Rosendal-Solecki} and the fact that $G$ has a comeager conjugacy class. By Lemma~\ref{nsi}, $H\cap G^\Sigma=G^\Sigma$.

We denote by $\Sym(\Sigma)$ the subgroup of $G\wr\sinf$ of elements of the form $((g_\omega),\sigma)$, where $g_\omega=e$ for all $\omega\in\Omega$ and $\supp(\sigma)\subset \Sigma$. Since $\Sym(\Sigma)\cong\sinf$ and the non-trivial normal subgroups of $\sinf$ are the finitary symmetric and alternating subgroups, which are of index $2^{\aleph_0}$, we know that for $H\cap \Sym(\Sigma)=\Sym(\Sigma)$. Now, $W'=G\wr_\Sigma\Sym(\Sigma)$ is generated by $G^\Sigma$ and $\Sym(\Sigma)$. Thus $H\geq W'$.

Following the proof of \cite[Th.\,1]{MR859950}, we choose a continuum of almost disjoint moieties $(\Sigma_\alpha)$ of $\Sigma$ and an involution $g_\alpha\in\sinf$ exchanging $\Sigma_\alpha$ with $\Omega\setminus\Sigma$ and fixing pointwise $\Sigma\setminus\Sigma_\alpha$. By the pigeonhole principle, for some $\alpha\neq\beta$, $g_\alpha$ and $g_\beta$ are in the same $H$-class and thus $g=g_\beta^{-1} g_\alpha$ is in $H$. Let $\Sigma'=g(\Sigma)$, then $G\wr_{\Sigma'}\Sym(\Sigma')=gW'g^{-1}\leq H$. Since $\Sigma\cup\Sigma'=\Omega\setminus g_\beta(\Sigma_\alpha\cap\Sigma_\beta)$ and $\Sigma\cap\Sigma'=\Omega\setminus g_\beta(\Sigma_\alpha\cup\Sigma_\beta)$ are infinite, the first lemma of \cite{MR859950} shows that $H_0=G\wr_{\Sigma\cup\Sigma'}\Sym(\Sigma\cup\Sigma')\subset H$.

Let us denote by $F$ the finite set $g_\beta(\Sigma_\alpha\cap\Sigma_\beta)$. For $f\in F$, let $G_f\leq G\wr_\Omega\sinf$ be the corresponding copy of $G$ acting on $\Lambda\times\{f\}$. The subgroup $H_f=H\cap G_f$ has small index and thus is open in $G_f$. Now
\[
H\geq\left(\Pi_{f\in F}H_f\right)\times H_0.
\]
The right-hand side being open (containing the pointwise stabilizer of a finite number of points), $H$ is open as well.
\end{proof}

% End literal retained input author-proof.tex

\begin{thebibliography}{99}
\bibitem{MR859950} J. D. Dixon, P. M. Neumann \& S. Thomas - “Subgroups of small index in infinite symmetric groups” , Bull. London Math. Soc. 18 (1986) no. 6, p. 580-586
\bibitem{MR1231710} W. Hodges, I. Hodkinson, D. Lascar \& S. Shelah - “The small index property for $\omega $-stable $\omega $-categorical structures and for the random graph” , J. London Math. Soc. (2) 48 (1993) no. 2, p. 204-218
\bibitem{Kechris-Rosendal} A. S. Kechris \& C. Rosendal - “Turbulence, amalgamation, and generic automorphisms of homogeneous structures” , Proc. London Math. Soc. (3) 94 (2007) no. 2, p. 302-350
\bibitem{Rosendal-Solecki} C. Rosendal \& S. Solecki - “Automatic continuity of homomorphisms and fixed points on metric compacta” , Israel J. Math. 162 (2007), p. 349-371
\end{thebibliography}
\end{document}
